\chapter{\label{chap:intro}Introdução}

\sigla{API}{Application Programming Interface}
\sigla{CRUD}{Create, Read, Update, Delete}
\sigla{HTTP}{Hypertext Transfer Protocol}
\sigla{IDE}{Integrated Development Environment}
\sigla{JPA}{Java Persistence API}
\sigla{JSON}{JavaScript Object Notation}
\sigla{PUCRS}{Pontifícia Universidade Católica do Rio Grande do Sul}
\sigla{REST}{Representational State Transfer}
\sigla{TCC}{Trabalho de Conclusão de Curso}
\sigla{UI}{User Interface}
\sigla{UX}{User Experience}

\abrev{VSCode}{Visual Studio Code}
\abrev{BD}{Banco de Dados}

\simbolo{$\pi$}{Constante com valor aproximado de 3{,}14}

Com o avanço acelerado da tecnologia e a popularização da internet, ferramentas digitais têm sido integradas de forma crescente ao cotidiano, impactando positivamente áreas como saúde, educação, mobilidade urbana e finanças pessoais. Essa digitalização viabiliza maior eficiência, automação de tarefas e acesso aprimorado à informação, promovendo mais controle e autonomia para o cidadão em suas tomadas de decisão.

No âmbito financeiro, observa-se o crescimento de plataformas digitais voltadas à gestão de recursos pessoais. O controle financeiro, que já foi visto como uma recomendação, é hoje uma necessidade premente diante do contexto econômico brasileiro. Em dezembro de 2024, o número de brasileiros endividados chegou a \textbf{73{,}51 milhões} \cite{SERASA2024}. Em janeiro de 2025, o percentual de famílias endividadas reduziu-se para \textbf{76{,}1\%}, queda de \textbf{0{,}6 p.p.} em relação a dezembro, embora o quadro permaneça preocupante \cite{PEIC2025}.

A saúde financeira dos brasileiros também apresentou leve melhora: o \textit{Índice de Saúde Financeira do Brasileiro} (I\hyp SFB), desenvolvido pela FEBRABAN em parceria com o Banco Central, subiu de \textbf{56{,}2} em 2023 para \textbf{56{,}7} em 2024 — o melhor resultado dos últimos três anos \cite{BCB2024,ISFB2024}. Complementando esse cenário, um levantamento indica que \textbf{cerca de 57 milhões} de brasileiros estão endividados sem perceber, e \textbf{19 milhões} já chegaram a negativar o nome, evidenciando a urgência de maior conscientização \cite{SERASA2025INFO}.

A má administração de recursos pode levar ao acúmulo de dívidas, perda de poder de compra e prejuízo à qualidade de vida. Por outro lado, uma gestão consciente viabiliza estabilidade, criação de reservas de segurança e o alcance de objetivos pessoais. A educação financeira, nesse sentido, emerge não apenas como ferramenta de empoderamento individual, mas também como base de decisões mais equilibradas e sustentáveis \cite{ANBIMA2024}.

Diante desse contexto, este trabalho apresenta o desenvolvimento e a implantação do \textbf{Dashboard Financeiro}, uma aplicação web voltada à gestão de finanças pessoais. Seu desenho \textit{full\hyp stack} organiza-se em duas camadas principais: o \textbf{frontend}, responsável pela interface e interação com o usuário; e o \textbf{backend}, responsável pela lógica de negócio, autenticação e persistência dos dados. Essa combinação modular favorece a criação de sistemas escaláveis e alinhados às boas práticas de engenharia de software.

O termo \textit{dashboard} refere-se a um painel visual que sintetiza dados e indicadores de forma clara e acessível, facilitando a identificação de padrões, o monitoramento de metas e a tomada de decisões embasadas \cite{MUNZNER2014,FEW2017}. No \textbf{Dashboard Financeiro}, o usuário registra transações, classifica receitas e despesas, define metas e acompanha visualizações que evidenciam comportamentos e tendências ao longo do tempo.

\section*{Objetivos}

\textbf{Objetivo Geral}

Desenvolver e implantar uma aplicação web \textit{full\hyp stack}, denominada \textbf{Dashboard Financeiro}, para organização e análise de finanças pessoais, utilizando práticas modernas, seguras, escaláveis e acessíveis de engenharia de software.

\textbf{Objetivos Específicos}

\begin{itemize}
    \item Permitir o registro e a categorização de receitas e despesas por meio de uma interface intuitiva;
    \item Apresentar gráficos e indicadores financeiros de forma dinâmica e responsiva;
    \item Oferecer recursos de definição de metas mensais e alertas de desempenho com notificações por e\hyp mail;
    \item Garantir segurança e privacidade dos dados por meio de autenticação baseada em tokens e criptografia;
    \item Estruturar a aplicação com arquitetura orientada a serviços e persistência em banco de dados relacional;
    \item Implantar a aplicação em ambiente de nuvem para acesso remoto e escalabilidade.
\end{itemize}

Todos os objetivos foram endereçados no volume atual: os capítulos~\ref{chap:chap3} e~\ref{chap:chap5} detalham a arquitetura orientada a serviços e a infraestrutura em produção, enquanto o Capítulo~\ref{chap:chap4} evidencia o cumprimento do cronograma e o Capítulo~\ref{chap:chap6} descreve a validação funcional e os próximos passos.

\section{Potencial de Impacto da Solução}

O \textbf{Dashboard Financeiro} visa transformar a forma como os usuários lidam com suas finanças, fomentando hábitos saudáveis e promovendo autonomia econômica. Sua abordagem educativa, visual e segura pode gerar impacto positivo na vida das pessoas — ajudando-as a planejar, economizar, evitar endividamentos desnecessários e tomar decisões mais conscientes, com efeitos benéficos para a sociedade como um todo.
